(1) Fourier分离的Laplace方程是 \(\Phi_{zz}-k^2\Phi=0\)，唯一Dirichlet解
\[
\Phi=V_1\frac{\sinh k(d-z)}{\sinh kd}\cos kx+
V_2\frac{\sinh kz}{\sinh kd}\cos(kx+\varphi).
\]
对 \(x,z\) 求负梯度得
\[
E_x=kV_1\frac{\sinh k(d-z)}{\sinh kd}\sin kx+
kV_2\frac{\sinh kz}{\sinh kd}\sin(kx+\varphi),
\]
\[
E_z=kV_1\frac{\cosh k(d-z)}{\sinh kd}\cos kx-
kV_2\frac{\cosh kz}{\sinh kd}\cos(kx+\varphi).
\]
在两端的 \(\Phi\) 恰好等于指定电位，且电位在间隙内满足 \(\nabla^2\Phi=0\)。

(2) 两端面向介质的法线分别为 \(+\mathbf e_z,-\mathbf e_z\)，故
\[
\sigma_0=\varepsilon k[V_1\coth(kd)\cos kx-
V_2\operatorname{csch}(kd)\cos(kx+\varphi)],
\]
\[
\sigma_d=\varepsilon k[V_2\coth(kd)\cos(kx+\varphi)-
V_1\operatorname{csch}(kd)\cos kx].
\]
用复幅度 \(v=(V_1,V_2e^{i\varphi})^\mathsf T\)，矩阵
\[
C_k=\varepsilon k
\begin{pmatrix}\coth kd&-\operatorname{csch}kd\\
-\operatorname{csch}kd&\coth kd\end{pmatrix}.
\]
本征值为 \(\varepsilon k(\coth kd\pm\operatorname{csch}kd)>0\)。
当 \(kd\gg1\)，\(C_k\to\varepsilon kI\)；当 \(kd\ll1\)，同相模本征值 \(\sim\varepsilon k^2d/2\)，反相模 \(\sim 2\varepsilon/d\)。这里的图案是由分段独立源施加的非均匀Dirichlet条件，不是一整块处处等势的金属板。

(3) 按 \(\langle\cos^2kx\rangle_x=1/2\)，利用Green恒等式，得到每单位面积场能
\[
\boxed{\bar{\mathcal U}
=\frac14\operatorname{Re}(v^\dagger C_kv)
=\frac{\varepsilon k}{4}
[(V_1^2+V_2^2)\coth(kd)-2V_1V_2\cos\varphi\,\operatorname{csch}(kd)]}.
\]
正定电容矩阵确保场能非负。
对于 \(k\to0^+\) 但始终对整周期 \(2\pi/k\) 平均，场能极限是
\[
\bar{\mathcal U}\to\frac{\varepsilon}{4d}
(V_1^2+V_2^2-2V_1V_2\cos\varphi).
\]
若先在有限 \(x\) 位置固定 \(k=0,\varphi=0\) 而把电极变成两个均匀电位，则其能量是 \(\varepsilon(V_1-V_2)^2/(2d)\)，恰为同相周期平均极限的两倍。原因是面积平均顺序改变，不能忽略平均 \(\cos^2\) 的不一致。

(4) Maxwell静电张量在二维的 \(zz\) 分量为
\(T_{zz}=\varepsilon(E_z^2-E_x^2)/2\)。
对 \(x\) 周期平均，\(\partial_z\langle T_{zz}\rangle=0\)，因为 \(\partial_iT_{iz}=0\) 且横向全周期导数积分为零。代入上下任一边界的场可得
\[
\langle T_{zz}\rangle
=\frac{\varepsilon k^2}{4}[
(V_1^2+V_2^2)\operatorname{csch}^2(kd)
-2V_1V_2\cos\varphi\,\operatorname{csch}(kd)\coth(kd)].
\]
另一方面，固定 \(V_1,V_2,\varphi,k\) 对第三问求间距导数，
\[
\frac{\partial\bar{\mathcal U}}{\partial d}
=-\frac{\varepsilon k^2}{4}[
(V_1^2+V_2^2)\operatorname{csch}^2(kd)
-2V_1V_2\cos\varphi\,\operatorname{csch}(kd)\coth(kd)].
\]
所以 \(\langle T_{zz}\rangle=-\partial_d\bar{\mathcal U}\)。
固定电势驱动的机械合力还须包含外电源所作功，不能仅对场能取负导数而遗漏源能量。
